%% 01-overview.tex — What is this portfolio?

\section{What Is This Portfolio?}

This portfolio is a collection of more than forty research papers asking a single question:
\emph{Can a computer draw fairer congressional districts than a human politician?}

The answer, across all fifty states and three census decades, is yes.

\subsection*{The Problem in One Sentence}

Every ten years, after the Census, each state redraws its congressional districts.
In most states, the party in power draws the lines---and draws them to stay in power.
This practice is called \textbf{gerrymandering}, and in 2019 the U.S.\ Supreme Court
ruled that federal courts can do nothing about it~\cite{rucho2019}.

\subsection*{The Proposal in One Sentence}

Replace the politician with an algorithm: a mathematical procedure that divides a
state into equal-population districts using only geography and population counts---no
voter registration, no election history, no racial or partisan data.

\subsection*{Five Research Tracks}

The portfolio is organized into seven tracks, labeled A through G,
each investigating a different aspect of algorithmic redistricting.
Together they span the question from first principles through
legal compliance, validation, and real-world application.

\begin{center}
\begin{tikzpicture}[
  track/.style={rectangle, draw=black, thick, rounded corners=4pt,
                minimum width=3.8cm, minimum height=1.1cm,
                align=center, font=\small\bfseries},
  arrow/.style={->, thick, >=stealth},
  node distance=0.45cm and 0.6cm
]

% Row 1 — foundation
\node[track, fill=blue!15]  (A) {Track A\\Overview \& Guides};
\node[track, fill=blue!25, right=of A]  (B) {Track B\\Algorithm};
\node[track, fill=green!20, right=of B] (C) {Track C\\Validation};

% Row 2 — applications
\node[track, fill=orange!20, below=0.7cm of A] (D) {Track D\\Legal / VRA};
\node[track, fill=red!15,    below=0.7cm of B] (E) {Track E\\Extensions};
\node[track, fill=purple!15, below=0.7cm of C] (F) {Track F\\State Leg.};

% Row 3 — synthesis
\node[track, fill=gray!20, below=0.7cm of E] (G) {Track G\\Ensemble};

% Arrows: B feeds everything
\draw[arrow] (B) -- (C);
\draw[arrow] (B) -- (D);
\draw[arrow] (B) -- (E);
\draw[arrow] (B) -- (F);
\draw[arrow] (B) -- (G);
\draw[arrow] (A) -- (B);

\end{tikzpicture}
\end{center}

\medskip

\noindent Table~\ref{tab:tracks} summarizes the seven tracks and their paper counts.

\begin{table}[h]
\centering
\caption{Portfolio tracks at a glance.}
\label{tab:tracks}
\begin{tabular}{@{}clll@{}}
\toprule
Track & Name & Core Question & Papers \\
\midrule
A & Overview \& Guides     & What does this portfolio do?        & 5 \\
B & Algorithm              & How does bisection work?            & 18 \\
C & Validation             & Does it hold up to scrutiny?        & 9 \\
D & Legal / VRA            & Is it lawful?                       & 6 \\
E & Extensions             & What else can it do?                & 7 \\
F & State Legislatures     & Does it scale to smaller chambers?  & 6 \\
G & Ensemble Analysis      & How does it compare to random maps? & 5 \\
\bottomrule
\end{tabular}
\end{table}
